\documentclass[10pt,a4paper]{amsart}
\usepackage[margin=19mm]{geometry}
\usepackage{amsmath,amssymb,mathtools}
\usepackage[scr=rsfs,cal=euler]{mathalfa}
\usepackage{xfrac}
\usepackage[unicode,hidelinks,bookmarks=false]{hyperref}
\newcommand{\ie}{i.e.\ }
\newcommand{\cf}{cf.\ }
\newcommand{\resp}{respectively\ }
\newcommand{\Subsection}{\subsection}
\providecommand{\nobreakdash}{\nobreak\hbox{-}}
\setlength{\emergencystretch}{2em}
\multlinegap0pt
\usepackage{mathtools}
\usepackage{xcolor}
\newtheorem{prop}{Proposition}
\title{On the generalized inverse tangent integral and Catalan’s constant}
\author{Petr Vlachopulos}
\date{Source: arXiv:2605.01830v1 (CC BY 4.0)}
\begin{document}
\maketitle

\noindent\emph{Source and licence.} On the generalized inverse tangent integral and Catalan’s constant. Version arXiv:2605.01830v1 (CC BY 4.0), distributed by arXiv under the Creative Commons Attribution 4.0 International licence, http://creativecommons.org/licenses/by/4.0/, as recorded in the arXiv licence metadata for this exact version and shown on the abstract page https://arxiv.org/abs/2605.01830v1. Full licence notice: \texttt{licenses/arxiv-2605.01830-CC-BY-4.0.txt}.\par\smallskip

\noindent\textit{Editorial scope.} Counted unit: the article's prop my.prop (source0.tex lines 537--548). The article's complete original proof of it is delivered with it (source0.tex lines 550--657, one \texttt{proof} environment of 108 lines). No mathematics is written, repaired or re-derived: the statement and the proof are the article's own lines, in the article's own notation, and no retained line is substituted or re-flowed.  No further source-local context is needed: every cross-reference of every retained line resolves inside the two delivered spans.  Nothing was omitted from any retained span, so the package carries no omission record.  No retained line contains a citation, so the wrapper carries no bibliography and no bibliography entry.  No retained line contains a figure environment, an image-inclusion command or drawing code, so no image is delivered and the package declares no asset dependency.  Editorial formatting only, in addition: an AMS \texttt{amsart} preamble that restates the article's own declarations and macros used by the retained lines, namely the environments \texttt{prop} on separate counters, together with the standard packages those lines need.  The retained environments print in this document as Proposition 1 in order of appearance, whereas the article prints the same items with the numbering of its own full document.  The complete native source of the article as distributed in the arXiv e-print for this exact version is bundled unchanged as \texttt{sources/199664/source0.tex}.
\begin{prop}
\phantomsection\label{my.prop}
For every \(a>0\), the inverse tangent integral satisfies
\begin{equation}
\operatorname{Ti}_2(a)
=
\arctan(a)\operatorname{Log}(a)
+
\mathscr{I}\!\bigl(\operatorname{Li}_2(1+ia)\bigr)
-\frac{\pi}{4}\operatorname{Log}(1+a^2).
\end{equation}
\end{prop}

\begin{proof}[Proof of Proposition \ref{my.prop}]\normalfont
Let us define the following quantity
\begin{equation}\label{eq.38}
H(a):=\arctan(a)\operatorname{Log}(a)
+
\mathscr{I}\!\bigl(\operatorname{Li}_2(1+ia)\bigr)
-\frac{\pi}{4}\operatorname{Log}(1+a^2),
\qquad a>0.
\end{equation}
We want to show that \(H'(a)=\arctan(a)/a\), which actually implies that
\begin{equation}
H(a)=\int_0^a \frac{\arctan x}{x}\,dx=\operatorname{Ti}_2(a),
\end{equation}
once the initial value is checked.

First, notice that
\begin{equation}\label{eq.40}
\frac{d}{da}\bigl(\arctan(a)\operatorname{Log}(a)\bigr)
=
\frac{\operatorname{Log}(a)}{1+a^2}+\frac{\arctan(a)}{a}.
\end{equation}

Next we compute the derivative of the dilogarithmic term in (\ref{eq.38}). Because \(1+ia\notin [1,\infty)\), then
for every \(a>0\), the derivative formula for \(\operatorname{Li}_2\) gives us
\begin{equation}
\frac{d}{da}\operatorname{Li}_2(1+ia)
=
-\frac{\operatorname{Log}(1-(1+ia))}{1+ia}\cdot i
=
-\frac{i\,\operatorname{Log}(-ia)}{1+ia}.
\end{equation}
For \(a>0\), the principal logarithm satisfies \(\operatorname{Log}(-ia)=\operatorname{Log}(a)-\frac{i\pi}{2}\). Therefore, we obtain the shift identity
\begin{equation}
-i\,\operatorname{Log}(-ia)
=
-i\operatorname{Log}(a)-\frac{\pi}{2}.
\end{equation}
Hence
\begin{equation}\label{eq.42}
\frac{d}{da}\operatorname{Li}_2(1+ia)
=
\frac{-\frac{\pi}{2}-i\operatorname{Log}(a)}{1+ia}.
\end{equation}
Multiply the numerator and the denominator of (\ref{eq.42}) by \(1-ia\) yields
\begin{equation}
\frac{d}{da}\operatorname{Li}_2(1+ia)
=
\frac{\bigl(-\frac{\pi}{2}-i\operatorname{Log}(a)\bigr)(1-ia)}{1+a^2}.
\end{equation}
Now expand the numerator in the following way
\begin{equation}\label{eq.44}
\bigl(-\tfrac{\pi}{2}-i\operatorname{Log}(a)\bigr)(1-ia)
=
-\frac{\pi}{2}
+
i\frac{a\pi}{2}
-
i\operatorname{Log}(a)
-
a\operatorname{Log}(a).
\end{equation}
Its imaginary part corresponds to \(\frac{a\pi}{2}-\operatorname{Log}(a)\).
Hence
\begin{equation}\label{eq.46}
\frac{d}{da}\mathscr{I}\!\bigl(\operatorname{Li}_2(1+ia)\bigr)
=
\mathscr{I}\!\left(\frac{d}{da}\operatorname{Li}_2(1+ia)\right)
=
\frac{\frac{a\pi}{2}-\operatorname{Log}(a)}{1+a^2}.
\end{equation}

Finally,
\begin{equation}\label{eq.47}
\frac{d}{da}\left(-\frac{\pi}{4}\operatorname{Log}(1+a^2)\right)
=
-\frac{\pi a}{2(1+a^2)}.
\end{equation}

Combining the relations (\ref{eq.40}), (\ref{eq.46}), and (\ref{eq.47}), we obtain
\begin{equation}
H'(a)
=
\left(\frac{\operatorname{Log}(a)}{1+a^2}+\frac{\arctan(a)}{a}\right)
+
\frac{\frac{a\pi}{2}-\operatorname{Log}(a)}{1+a^2}
-
\frac{\pi a}{2(1+a^2)}
=
\frac{\arctan(a)}{a}.
\end{equation}
Therefore
\begin{equation}
H(a)=\int_0^a \frac{\arctan x}{x}\,dx + C
\end{equation}
for some integration constant \(C\in\mathbb{R}\).

In order to determine the integration constant \(C\), we should let \(a\to 0^+\) in (\ref{eq.38}). We get
\[
\arctan(a)\operatorname{Log}(a)\to 0,
\qquad
\mathscr{I}(\operatorname{Li}_2(1+ia))\to \mathscr{I}(\operatorname{Li}_2(1))=0,
\qquad
\operatorname{Log}(1+a^2)\to 0.
\]
Hence \(H(a)\to 0\), which actually implies that \(\operatorname{Ti}_2(0)=0\). Thus \(C=0\), and the claimed formula naturally follows.

This completes the proof of Proposition \ref{my.prop}.
\end{proof}

\end{document}
